% ============================================================
% Page 8 — Alban (au XIIème siècle : époque cathare)
% Transcription du document source.
% Guillemets fermants ajoutés à la demande de l'utilisateur.
% Texte conservé sans autre reformulation ni correction.
% ============================================================

\thispagestyle{plain}

\vspace*{12mm}

\begin{center}
    {\Large\bfseries Alban (au XII\textsuperscript{ème} siècle : époque cathare)}
\end{center}

\vspace{18mm}

\noindent
Un moine -- Henri de Lausanne --, un prêtre provençal -- Pierre de
Bruys -- font à Albi le lit à l'hérésie albigeoise. « Le premier
prêche, un christianisme sans le Christ, sans culte, sans sacrements,
sans prêtre, sans église. Le mariage est déclaré « institution
diabolique ». Le second brûle les crucifix, les statues, sous prétexte
d'épuration du culte -- de là naît une décomposition sociale que le
dualisme cathare accentuera. »

\par\vspace{12mm}

\noindent
Le 29 juin 1145, Saint Bernard tonne à la cathédrale d'Albi contre
le dévergondage des mœurs et une piété déliquescente. « Il n'y a
plus que des Chrétiens sans Christ », s'écrie le moine blanc. Ses
invectives n'épargnent pas le clergé. « Les Apôtres étaient la lumière
du monde, vous en êtes les ténèbres. Vous traînez des femmes avec
vous, voire même vous les enfermez avec vous. Comment votre vertu
est-elle en sûreté dans une telle compagnie ? Etre toujours avec
une femme et ne pas pécher est plus difficile que de ressusciter
un mort. »

\par\vspace{12mm}

\noindent
Albi se convertit et le foyer de l'hérésie se déplace pour se fixer
dans la région de Lombers. Qu'advint-il à Alban ? Fut-il contaminé ?
On sait que le vicomte d'Ambialet soutint les Albigeois. Sa vicomté
comprenait Ginestous, La Balme... aux portes d'Alban. Rien ne permet,
à ce jour, d'avancer une certitude à ce sujet.

\vfill
